\section{切空间与维数}\label{sec:dimension-of-tangent-spaces}

在第~\ref{sec:linearizations-and-tangent-spaces}~节看到线性化和切空间的几个例子. 读者或许注意到, 切空间的维数有时大于簇本身的维数, 却从不小于它. 不可约簇的维数是该簇任一点处切空间维数的下界, 这就是本节的主要结论.

\begin{Proposition}{切空间维数的下界}\leavevmode\label{prop:tangent-dimension-lower-bound}
设 $X\subseteq\mathbb{A}^{n}$ 为不可约仿射簇, 且 $a\in X$. 则
\[
  \dim(T_{a}X)\geq\dim(X).
\]
\end{Proposition}

注意, 不等式左边的维数是 $T_{a}X\subseteq K^{n}$ 作为向量空间的维数, 右边的维数是 $X\subseteq\mathbb{A}^{n}$ 作为仿射簇的维数. 如果愿意, 可以将 Proposition~\ref{prop:tangent-dimension-lower-bound} 中不等式的两边都解释为 $\mathbb{A}^{n}$ 中仿射簇的维数, 只需注意 $\dim(T_{a}X)=\dim(L_{a}X)$ (Exercise~\ref{ex:tangent-space-and-linearization-have-equal-dimension}).

在证明 Proposition~\ref{prop:tangent-dimension-lower-bound} 之前, 先证明零维情形中一个更强的充要条件.

\begin{Lemma}{零维切空间}\leavevmode\label{lem:zero-dimensional-tangent-spaces}
若 $X\subseteq\mathbb{A}^{n}$ 为不可约仿射簇, 且 $a\in X$, 则 $\dim(X)=0$ 当且仅当 $\dim(T_{a}X)=0$.
\end{Lemma}
\begin{Proof}
若 $\dim(X)=0$, 则 $X$ 只含点 $a$. 由定义可验证 $T_{a}X=\{0\}$ (Exercise~\ref{ex:tangent-space-of-single-point}).

为证明反向, 假设 $\dim(T_{a}X)=0$, 并考虑由在 $a$ 处取值为零的多项式函数组成的极大理想 $I_{a}\subseteq K[X]$. 我们要证明 $I_{a}$ 是零理想. 这将推出 $K[X]$ 中的零理想是极大理想, 所以 $K[X]$ 是域, 再由零点定理推出 $X$ 只有一个点.

因为有限维向量空间与其对偶维数相同, Theorem~\ref{thm:point-ideal-square-quotient-is-dual-of-tangent-space} 与 $\dim(T_{a}X)=0$ 这一假设蕴含 $I_{a}=I_{a}^{2}$. 选取生成元
\[
  I_{a}=\langle F_{1},\ldots,F_{m}\rangle.
\]
由于 $I_{a}=I_{a}^{2}$, 可以把每个 $F_{i}$ 视为 $I_{a}^{2}$ 的元素, 从而可以找到形如
\[
  F_{i}=G_{i,1}F_{1}+\cdots+G_{i,m}F_{m}
\]
的等式, 其中对每个 $i,j=1,\ldots,m$, 有 $G_{i,j}\in I_{a}$. 于是得到线性方程组
\[
  (I-G)\cdot
  \begin{pmatrix}F_{1}\\\vdots\\F_{m}\end{pmatrix}
  =\begin{pmatrix}0\\\vdots\\0\end{pmatrix},
\]
其中 $I$ 为 $m\times m$ 单位矩阵, $G$ 是以 $G_{i,j}$ 为元素的矩阵.

应用 Cramer 法则, 得对所有 $i=1,\ldots,m$ 都有
\[
  \det(I-G)F_{i}=0\in K[X].
\]
注意 $\det(I-G)$ 是 $X$ 上的多项式函数, 又因为对所有 $i,j$ 都有 $G_{i,j}(a)=0$, 所以 $\det(I-G)$ 在 $a$ 处取值为 $1$. 因此 $\det(I-G)$ 不是 $X$ 上的零函数. 因为 $K[X]$ 是整环 (这里用到 $X$ 不可约的假设), 可知对所有 $i$ 都有 $F_{i}=0$. 所以 $I_{a}$ 是零理想, 证明完毕.
\end{Proof}

现在证明 Proposition~\ref{prop:tangent-dimension-lower-bound}. 证明对 $\dim(X)$ 归纳. 归纳步骤 (需要用到 Lemma~\ref{lem:zero-dimensional-tangent-spaces}) 通过用一个超平面, 即由单个线性方程定义的簇, 同时切割 $X$ 与 $L_{a}X$, 并应用维数论基本定理来完成.

\begin{Proof}\textbf{of Proposition~\ref{prop:tangent-dimension-lower-bound}.}\quad
给定不可约仿射簇 $X\subseteq\mathbb{A}^{n}$ 及点 $a\in X$, 对 $\dim(X)$ 归纳证明 $\dim(T_{a}X)\geq\dim(X)$.

(归纳起点) 若 $\dim(X)=0$, 则由向量空间的维数非负可知 $\dim(T_{a}X)\geq\dim(X)$. (实际上, 由 Lemma~\ref{lem:zero-dimensional-tangent-spaces} 还知道更多: 在这种情形下 $\dim(T_{a}X)=\dim(X)$.)

(归纳步骤) 设 $X\subseteq\mathbb{A}^{n}$ 为正维不可约仿射簇. 假设命题中的不等式对所有维数为 $\dim(X)-1$ 的不可约仿射簇成立. 令 $a=(a_{1},\ldots,a_{n})\in X$.

由于已假设 $\dim(X)>0$, Lemma~\ref{lem:zero-dimensional-tangent-spaces} 中\textquotedbl{}若\textquotedbl{}的方向蕴含 $T_{a}X\supsetneq\{0\}$, 这等价于 $L_{a}X\supsetneq\{a\}$. 因此可以选取一点 $b\in L_{a}X\setminus\{a\}$. 两点 $a,b\in\mathbb{A}^{n}$ 至少有一个坐标不同, 不妨设它们的第一个坐标不同, 并定义超平面
\[
  H=\mathcal{V}(x_{1}-a_{1})\subseteq\mathbb{A}^{n}.
\]
令 $Y=X\cap H$. 首先注意 $Y$ 不可能是整个 $X$. 事实上, 若 $Y=X$, 则 $X\subseteq H$, 这将蕴含 $L_{a}X\subseteq H$ (Exercise~\ref{ex:linearization-monotone-under-inclusion} 和~\ref{ex:linear-variety-equals-its-linearization}). 但选 $H$ 时特意使得 $b\in L_{a}X$ 而 $b\notin H$, 所以 $L_{a}X\not\subseteq H$. 因此 $Y\neq X$. 又由于 $a\in Y$, 可知
\begin{equation}\label{eq:tangent-lower-bound-proper-slice}
  \varnothing\subsetneq Y\subsetneq X,
\end{equation}
故由维数论基本定理的强版本, $Y$ 的每个不可约分支都有维数 $\dim(X)-1$.

令 $Z$ 为 $Y$ 的一个包含 $a$ 的不可约分支. 由于 $Z\subseteq X$ 且 $Z\subseteq H$, 有 $L_{a}Z\subseteq L_{a}X$ (Exercise~\ref{ex:linearization-monotone-under-inclusion}) 和 $L_{a}Z\subseteq H$ (Exercise~\ref{ex:linear-variety-equals-its-linearization}). 这蕴含
\begin{equation}\label{eq:tangent-lower-bound-linearized-slice}
  L_{a}Z\subseteq L_{a}X\cap H\subsetneq L_{a}X,
\end{equation}
从而 $T_{a}Z\subsetneq T_{a}X$, 所以 $\dim(T_{a}Z)<\dim(T_{a}X)$.

综上, 并对 $Z$ 用归纳假设, 得
\[
  \dim(X)-1=\dim(Z)\leq\dim(T_{a}Z)\leq\dim(T_{a}X)-1,
\]
故 $\dim(X)\leq\dim(T_{a}X)$, 证明完毕.
\end{Proof}

虽然 Proposition~\ref{prop:tangent-dimension-lower-bound} 给出了 $T_{a}X$ 维数的下界, 但不存在只用 $\dim(X)$ 表示的上界. 为说明这一点, 下面的例子给出构造仿射曲线的一种方法, 使曲线上有一点的切空间维数可以任意大.

\begin{Example}{小簇的大切空间}\leavevmode\label{ex:curves-with-arbitrarily-large-tangent-spaces}
固定 $n\geq1$, 考虑仿射簇
\[
  X_{n}=\mathcal{V}\bigl(\{x_{i}^{n+j-1}-x_{j}^{n+i-1}\mid1\leq i<j\leq n\}\bigr)
    \subseteq\mathbb{A}^{n}.
\]
注意 $X_{2}$ 是 Example~\ref{ex:linearizations-of-cusp} 中的尖点曲线. 可以证明 (Exercise~\ref{ex:monomial-affine-curve-parametrization-and-dimension}), $X_{n}$ 是一个一维不可约仿射簇, 它的点形如
\[
  X_{n}=\{(b^{n},b^{n+1},\ldots,b^{2n-1})\mid b\in\mathbb{A}^{1}\}.
\]
令 $a=(0,\ldots,0)\in X$. 我们论证 $L_{a}X_{n}=\mathbb{A}^{n}$, 因而 $T_{a}X_{n}=K^{n}$.

回忆 $L_{a}X_{n}=\mathcal{V}(\{L_{a}f\mid f\in\mathcal{I}(X)\})$. 考虑多项式
\[
  f\in\mathcal{I}(X_{n})\subseteq K[x_{1},\ldots,x_{n}].
\]
由于 $f$ 在 $a=(0,\ldots,0)\in X_{n}$ 处取值为零, 它的常数项为零, 所以可以写成
\[
  f=\sum_{i=1}^{n}c_{i}x_{i}+g,
\]
其中 $c_{i}\in K$, 且 $g$ 的常数项和线性项全为零. 在点 $(b^{n},\ldots,b^{2n-1})\in X_{n}$ 处取值, 有
\[
  0=f(b^{n},\ldots,b^{2n-1})
  =\sum_{i=1}^{n}c_{i}b^{n+i-1}+g(b^{n},\ldots,b^{2n-1}).
\]
因为这个关系对所有 $b\in\mathbb{A}^{1}$ 成立, 相应的多项式是零多项式:
\[
  0=\sum_{i=1}^{n}c_{i}y^{n+i-1}+g(y^{n},\ldots,y^{2n-1})\in K[y].
\]
由于 $\sum_{i=1}^{n}c_{i}y^{n+i-1}$ 中各项的 $y$ 的次数各不相同且均小于 $2n$, 而 $g(y^{n},\ldots,y^{2n-1})$ 中各项的次数至少为 $2n$, 因此对所有 $i$ 都有 $c_{i}=0$. 所以 $f$ 的线性项为零, 即 $L_{a}f=0$. 因为这对所有 $f\in\mathcal{I}(X_{n})$ 成立, 得到 $L_{a}X_{n}=\mathcal{V}(0)=\mathbb{A}^{n}$ 和 $T_{a}X_{n}=K^{n}$.
\end{Example}

\BookPairedNote
  {前例有一个关于仿射簇嵌入的有趣推论, 可以与光滑流形理论中的 Whitney 嵌入定理对照. 由于 $X_{n}$ 在原点处的切空间是 $n$ 维的, 任何与 $X_{n}$ 同构的仿射簇也会有一个 $n$ 维的切空间.}
  {Whitney 嵌入定理断言任意 $n$ 维光滑流形可视为 $\mathbb{R}^{2n}$ 的子流形.}
特别地, 曲线 $X_{n}$ 不可能与 $\mathbb{A}^{m}$ 内的任何仿射簇同构, 只要 $m<n$. 因此对任意 $m>0$, 这个例子表明存在不能嵌入 $\mathbb{A}^{m}$ 的一维仿射簇.

\BookExerciseHeading
\begin{BookExercises}
  \item\label{ex:tangent-space-of-single-point} 若 $X=\{a\}\subseteq\mathbb{A}^{n}$ 只有一个点, 证明 $T_{a}X=\{0\}\subseteq K^{n}$.
  \item\label{ex:linearization-monotone-under-inclusion} 设 $X\subseteq Y\subseteq\mathbb{A}^{n}$ 为仿射簇. 对任意 $a\in X$, 证明
  \[
    L_{a}X\subseteq L_{a}Y.
  \]
  由此推出 $T_{a}X\subseteq T_{a}Y$.
  \item\label{ex:linear-variety-equals-its-linearization} 假设 $X=\mathcal{V}(\ell_{1},\ldots,\ell_{m})\subseteq\mathbb{A}^{n}$, 其中 $\ell_{1},\ldots,\ell_{m}\in K[x_{1},\ldots,x_{n}]$ 都是线性多项式. 证明对任意 $a\in X$, 都有 $L_{a}X=X$.
  \item\label{ex:monomial-affine-curve-parametrization-and-dimension} 设 $X_{n}$ 为 Example~\ref{ex:curves-with-arbitrarily-large-tangent-spaces} 中定义的仿射簇.
  \begin{enumerate}
    \item 证明
    \[
      X_{n}=\{(b^{n},b^{n+1},\ldots,b^{2n-1})\mid b\in\mathbb{A}^{1}\}.
    \]
    (提示: 取 $b$ 为前两个坐标的商.)
    \item 证明 $\dim(X_{n})=1$. (提示: 证明 $\{[x_{1}]\}$ 是一个 Noether 基.)
    \item 证明 $X_{n}$ 不可约. (提示: (a) 可能有用.)
  \end{enumerate}
  \item\label{ex:isolated-real-point-with-two-dimensional-tangent-space} 本题说明切空间在实数域上会有一些奇特的表现. 考虑不可约实仿射簇
  \[
    X=\mathcal{V}(y^{2}-x(x+1)^{2})\subseteq\mathbb{A}^{2}_{\mathbb{R}}.
  \]
  \begin{enumerate}
    \item 证明 $(-1,0)$ 为 $X$ 的孤立点. 换句话说, 证明 $(-1,0)\in X$, 且能找到以 $(-1,0)$ 为中心, 半径为某个正数的圆盘 $D$, 使得 $D\cap X=\{(-1,0)\}$.
    \item 由于 $(-1,0)$ 是 $X$ 的孤立点, 你可能预期 $\dim(T_{(-1,0)}X)=0$. 与此相反, 证明 $\dim(T_{(-1,0)}X)=2$.
    \item 如前面的例子所见, 多项式方程的实解通常不能显示全貌, 往往还需要看复解. 若考虑复簇
    \[
      X_{\mathbb{C}}=\mathcal{V}(y^{2}-x(x+1)^{2})\subseteq\mathbb{A}^{2}_{\mathbb{C}},
    \]
    那么你认为在点 $(-1,0)\in X_{\mathbb{C}}$ 处发生了什么, 能帮助解释 (b) 的答案?
  \end{enumerate}
\end{BookExercises}
